% !TEX root =  ../Dissertation.tex

\chapter{Introduction}
In October 2021, Wordle was released and was a resounding success almost immediately after launch. The website grew traction quickly, going from only 90 players on November 1st to 300,000 on January 2nd 2022.\cite{Daniel2022} This success led to other developers working on games inspired by Wordle and has created a thriving community of smaller websites \cite{Adory2023} that follow in what will be referred to as the daily game genre of web application. The daily game genre has had continued success, but only a fraction of what Wordle has received.\cite{Google2024} Josh Wardle, the creator of Wordle, attributes much of the success of Wordle to the outreach of its shareable ‘emoji grid’  that was popular on Twitter early on in Wordle's conception.\cite{Slate2022} This was pivotal to bringing users to the website (which has since been acquired by NYT Games and added to their selection).\cite{nyt2022} However, once they gained this considerable player base, there was a lack of functionality that made players want to continue; the only thing that drew them back to the web application was the game itself. 

This project hypothesises that if retention techniques found in other multiplayer online games are added to the daily game genre of web applications \cite{Katkoff2016}, users will be more satisfied and retention will be increased. If the hypothesis is valid, developers in the daily game genre could augment their offerings to deliver a more engaging and exciting experience whilst increasing user retention.

To test this paper's hypothesis, there will be a self-conducted study of the effectiveness of these retention techniques by creating a simple Wordle clone, which will be referred to as Wordle no-login and a version of Wordle that has the retention techniques of more features, a leaderboard, and a set of achievements for participants to play, that will be named The Gauntlet. Users will be asked to complete two surveys, and with their permission, the web application will track the activity on their web browser to see if there is any increase in time spent and thus, engagement. The users will be asked to complete a survey when the testing is over to determine if they felt more satisfied using the website with increased functionality.